\documentclass{article}

% Un-anonymized preprint for course submission.
\usepackage[preprint]{neurips_2026}

\usepackage[utf8]{inputenc}
\usepackage[T1]{fontenc}
\usepackage{hyperref}
\usepackage{url}
\usepackage{booktabs}
\usepackage{amsfonts}
\usepackage{amsmath}
\usepackage{nicefrac}
\usepackage{microtype}
\usepackage{xcolor}
\usepackage{listings}

\lstset{
  basicstyle=\ttfamily\footnotesize,
  breaklines=true,
  columns=fullflexible,
  keepspaces=true,
  frame=single,
  framesep=3pt,
  xleftmargin=6pt,
  xrightmargin=6pt,
}

\title{A Dual-LLM System with Machine-Learned Cascade Routing and
Bounded Conversational Memory}

\author{%
  Zichen Qi \\
  Khoury College of Computer Sciences\\
  Northeastern University\\
  Boston, MA 02115 \\
  \And
  Yalin Sun \\
  Khoury College of Computer Sciences\\
  Northeastern University\\
  Boston, MA 02115 \\
  \And
  Sandeep Vijayarao \\
  Khoury College of Computer Sciences\\
  Northeastern University\\
  Boston, MA 02115 \\
  \texttt{vijayarao.s@northeastern.edu} \\
}

\begin{document}

\maketitle

\begin{abstract}
Large language models (LLMs) deliver strong reasoning quality, but routing every
user query through a frontier model is wasteful for the short, factual, or
conversational queries that dominate real traffic. We present a
\emph{dual-LLM system} that pairs a small model (GPT-4o-mini) with a large
model (GPT-4o) under a two-stage \emph{cascade router}: a fast offline classifier
(TF-IDF features $+$ hand-engineered tags $+$ logistic regression) answers on
its own when confident, and falls back to an LLM classifier only when it is
not. For ``big'' queries, the large model produces a topic-neutral core
answer, a one-line compression is computed, and the small model adds a
personalized intro and closing conditioned on the user profile and the
compressed core, so the large model's full output never re-enters the
small-model context. To keep long conversations tractable, we introduce a
\emph{bounded conversational buffer} that keeps the last $N$ turn-pairs
verbatim and folds older turns into an incremental rolling summary produced
by the small model in $\mathcal{O}(1)$ per turn. On a held-out split of our
hand-labelled seed set the offline router reaches
$95.5\%$ accuracy, the end-to-end system passes $38/38$ pytest integration
tests, and the ML router eliminates the LLM classifier call on high-confidence
queries, removing one API round-trip from the small-query path entirely.
Code is available at \url{https://github.com/sandeepvijayarao09/dual-llm-system}.
\end{abstract}

\section{Introduction}
\label{sec:intro}

Consumer LLM assistants face a cost--quality trade-off on every turn:
frontier-class models such as GPT-4o, Claude Opus, and Gemini Ultra deliver
the highest answer quality but cost one to two orders of magnitude more per
token than lighter-weight models such as GPT-4o-mini or Claude Haiku
\citep{openai_pricing_2024}. Crucially, a large fraction of conversational
traffic---greetings, short factual lookups, unit conversions,
quick code edits---does not need a frontier model at all. Recent work on
\emph{LLM cascades} and \emph{routers} \citep{chen_frugalgpt_2023,
ong_routellm_2024} shows that, on public benchmarks, 40--70\% of queries
can be handled by a much cheaper model with no measurable loss in quality,
provided we can decide \emph{per query} which model to use.

Independently, long multi-turn conversations pose a second scaling problem:
passing the entire history into every call makes the context window grow
linearly in turn count, inflates latency, and eventually truncates the
earliest turns \emph{without} warning. Production frameworks such as
LangChain have converged on \emph{summary-buffer memory}
\citep{langchain_memory_2023}---keep the last few turns verbatim and roll the
rest into a compact running summary. This is inexpensive if the summarizer is
itself a small model.

This paper describes a working system that combines both ideas into a single
orchestrator. Our contributions are:

\begin{enumerate}
    \item A \textbf{two-stage cascade router} in which an offline
    TF-IDF\,$+$\,logistic-regression classifier (the ``ML router'') answers on
    its own when its confidence exceeds a threshold $\tau$, and otherwise
    defers to an LLM classifier. This removes one small-LLM round-trip from
    every confident query.
    \item A \textbf{personalization-by-compression} dispatch for ``big'' queries:
    the large model produces a profile-free core answer; a one-line
    compression is computed by the small model; and the small model then
    generates a personalized intro/closing conditioned \emph{only} on that
    compression plus the user profile. The large model's full output never
    re-enters the small-model context, bounding its prompt size.
    \item A \textbf{bounded conversational buffer} that keeps the last $N$
    turn-pairs verbatim and incrementally folds older pairs into a rolling
    summary in $\mathcal{O}(1)$ per turn, with graceful degradation if the
    small LLM is unavailable.
    \item A \textbf{routing log} that stores every (query, ML decision,
    LLM decision, final decision) tuple in SQLite, so the router can be
    continuously retrained from production traffic, with a one-line CLI for
    exporting an agreement-only CSV training set.
    \item An \textbf{end-to-end open-source implementation} in Python with a
    pytest suite covering router features, model training, the conversation
    buffer, the classification logger, and the full orchestrator cascade
    under mocked LLMs (38 tests; all passing).
\end{enumerate}

\section{Related Work}
\label{sec:related}

\paragraph{LLM cascades and routers.}
\citet{chen_frugalgpt_2023} introduced FrugalGPT, which sequentially queries
progressively stronger LLMs and stops once a learned scoring function is
confident in the answer. \citet{ong_routellm_2024} train a pairwise router
that decides \emph{before} the call which of two models to use, demonstrating
large cost reductions with minimal quality loss on MT-Bench. Commercial
systems such as \citet{databricks_routing_2024} follow the same
predict-then-route pattern. Our router is in this second family---we decide
before the call---but adds a cascade inside the routing stage itself:
cheap offline classifier first, LLM classifier only on low-confidence inputs.

\paragraph{Conversational memory.}
The Transformer architecture \citep{vaswani_attention_2017} imposes a fixed
context window, which in a naive chat loop grows linearly with turn count.
Practical systems typically use either a hard sliding window (drop old turns)
or a summary buffer \citep{langchain_memory_2023, schlag_rmt_2023} in which
evicted turns are compressed. Our ConversationBuffer is a direct descendant
of LangChain's \texttt{ConversationSummaryBufferMemory}, but differs in two
ways: summarization is \emph{incremental} (one call per eviction rather than
a full re-summarization), and the buffer degrades gracefully to a
length-tagged textual trail when no small-LLM client is wired up, which makes
unit tests fully offline.

\paragraph{Personalization in LLM assistants.}
Prior work on persona-grounded generation
\citep{zhang_personachat_2018, salemi_lamp_2024} shows that lightweight user
profiles consistently improve subjective helpfulness. We implement a
minimalist form of this: a JSON profile per user, updated once per session
from the small LLM's inferences, and injected only into the small-LLM
prompts---never into the large-LLM prompt. This keeps profile PII out of the
frontier-model call path and is consistent with cost--privacy guidance in
recent deployment surveys \citep{bommasani_foundation_2021}.

\section{System Overview}
\label{sec:system}

Figure-style description: the orchestrator \citep[][\texttt{orchestrator.py}]{our_repo}
exposes a single \texttt{process(query, user\_id)} entry point. On every
call it (1) loads the user profile from SQLite, (2) retrieves the user's
ConversationBuffer, (3) applies a pre-flight length check that auto-routes
very long inputs to the large model, (4) runs the cascade router, (5)
dispatches to the small or large pipeline, (6) appends the turn to the
buffer, and (7) logs the decision to a SQLite routing log.

\begin{table}[h]
  \caption{System components and their responsibilities. ``S''\,=\,small LLM,
  ``L''\,=\,large LLM, ``OFFLINE''\,=\,no network I/O.}
  \label{tab:components}
  \centering
  \begin{tabular}{lll}
    \toprule
    Component & Calls & Role \\
    \midrule
    MLRouter       & OFFLINE & First-stage offline classifier (sklearn) \\
    Classifier     & S       & Fallback LLM classifier (JSON output) \\
    SimpleAnswerer & S       & Direct answer on the small route \\
    Reasoner       & L       & Core answer on the big route \\
    Personalizer   & S       & Summarize $+$ add intro/closing \\
    ProfileUpdater & S       & Post-session profile update \\
    ConversationBuffer & S (rare) & Bounded history + incremental summary \\
    ClassificationLogger & OFFLINE & SQLite log of every routing decision \\
    \bottomrule
  \end{tabular}
\end{table}

\section{Method}
\label{sec:method}

\subsection{The cascade router}
\label{sec:router}

Given a query $q$, the router outputs one of two labels,
$y \in \{\mathrm{small}, \mathrm{big}\}$. Let $f_{\mathrm{ML}}(q)
= (\hat{y}_{\mathrm{ML}}, c_{\mathrm{ML}})$ be the ML router, returning a
label and a confidence score $c_{\mathrm{ML}} \in [0,1]$, and let
$f_{\mathrm{LLM}}(q)$ be the LLM classifier. Our dispatch rule is:
\begin{equation}
  \mathrm{route}(q) \;=\;
  \begin{cases}
    \hat{y}_{\mathrm{ML}} & \text{if } c_{\mathrm{ML}} \geq \tau, \\
    f_{\mathrm{LLM}}(q) & \text{otherwise},
  \end{cases}
  \label{eq:cascade}
\end{equation}
with $\tau = 0.65$ in our configuration. This is the \emph{minimum} dispatch
rule that (i) never \emph{increases} latency (every ML call is cheap and can
be skipped by setting $\tau{=}0$), (ii) is a strict improvement when
$c_{\mathrm{ML}} \geq \tau$ has high precision, and (iii) preserves the LLM
classifier's behaviour on ambiguous inputs.

\paragraph{Features.} Rather than training a character-level or embedding-based
model, we use a simple hybrid feature set in which hand-engineered tags are
textually prepended to the query before it is fed into a standard TF-IDF
vectorizer. The tags (implemented in \texttt{router/features.py}) encode:
\begin{itemize}
    \item length buckets \texttt{LEN\_SHORT/MEDIUM/LONG};
    \item reasoning keywords \texttt{HAS\_REASONING\_KW} (\emph{why, how, compare,
    prove, explain, analyze, derive, justify, contrast, evaluate, reason, design,
    architect, optimize, debug});
    \item factual keywords \texttt{HAS\_FACTUAL\_KW} (\emph{what is, who is,
    when, where, capital, define, \ldots});
    \item greeting/thanks keywords \texttt{HAS\_GREETING\_KW};
    \item code or math presence (\texttt{HAS\_CODE}, \texttt{HAS\_MATH});
    \item question-mark density (\texttt{Q\_SINGLE} vs.\ \texttt{Q\_MULTI});
    \item sentence count (\texttt{MULTI\_SENTENCE}).
\end{itemize}
Tags are then concatenated with the raw query and the whole string is fed
into a sklearn \texttt{Pipeline}:
\begin{lstlisting}[language=Python]
Pipeline([
    ("tfidf", TfidfVectorizer(ngram_range=(1,2),
                              min_df=1, max_features=3000,
                              sublinear_tf=True)),
    ("clf",   LogisticRegression(C=1.0,
                                 class_weight="balanced",
                                 max_iter=1000, solver="liblinear")),
])
\end{lstlisting}
Concatenating tag tokens with the raw query makes them visible to both the
unigram and bigram features without requiring a custom featurizer, and keeps
the whole pipeline trivially serializable with \texttt{joblib}.

\paragraph{Seed data and training.} We hand-labelled a seed of 108 queries
spanning greetings, short factual lookups, unit conversions,
proofs, system design, research synthesis, and code debugging
(\texttt{router/seed\_data.py}). Training (\texttt{router/train.py}) does an
80/20 stratified split, fits the pipeline, reports metrics, then refits on
the full dataset and serializes to
\texttt{router/models/router\_v0.joblib}. The trainer also accepts an
\texttt{--extra} CSV argument, which is intended to ingest the
agreement-only export from the routing log; this is the continuous-learning
path (see \S\ref{sec:log}).

\subsection{Personalization by compression}
\label{sec:compress}

On the ``big'' route we do not simply append the large model's output to a
personalization prompt, because that would (a) force the large model's
full output through the small model's context on \emph{every} subsequent
turn (via the conversation buffer), and (b) expose any user PII in the
personalization prompt to the frontier-model call. Instead we do:
\begin{enumerate}
    \item $a = \mathrm{Reasoner}(q)$ \hfill (large LLM, no profile)
    \item $s = \mathrm{Summarize}(a)$ \hfill (small LLM, one sentence)
    \item $(i, c) = \mathrm{AddContext}(q, \mathrm{profile}, s)$ \hfill (small LLM)
    \item $\mathrm{final} = i \,\Vert\, a \,\Vert\, c$
\end{enumerate}
The small model sees only the one-sentence summary $s$, not the full $a$, so
the personalization prompt stays bounded regardless of the length of $a$.
The user profile is never forwarded to the large LLM.

\subsection{The conversation buffer}
\label{sec:buffer}

\texttt{ConversationBuffer} (\texttt{memory/conversation\_buffer.py}) maintains
a deque of recent messages and a string \texttt{rolling\_summary}. On each
\texttt{add\_turn(user, asst)} it appends two messages (the user turn and
the assistant turn). If the deque exceeds $2N$ messages (where $N$ is the
configurable window size in \emph{turn-pairs}, default 3), the oldest
turn-pair is popped and folded into the rolling summary:
\begin{equation}
  \mathrm{summary}_{t+1} \;=\; \mathrm{SmallLLM}\bigl(
    \mathrm{summary}_t, \;(\mathrm{user}_\mathrm{old}, \mathrm{asst}_\mathrm{old})
  \bigr),
  \label{eq:summary}
\end{equation}
where the small LLM is prompted to incorporate the new turn-pair into the
existing summary under an 80-token budget. Assistant turns longer than
$M = 400$ characters are truncated before being fed to the summarizer to
bound the input size of the eviction call. This makes the amortized cost of
\texttt{add\_turn} exactly one small-LLM call per eviction, i.e.\ $\mathcal{O}(1)$
per turn. If no small-LLM client is wired (e.g.\ in unit tests), the buffer
falls back to a textual concatenation of evicted turns so the calling code
remains correct.

\texttt{build\_history()} returns the list of chat messages consumed by the
LLM call: if the rolling summary is non-empty, it is prepended as a
\texttt{system} message; the remaining verbatim turns follow in order.
Per-user isolation is implemented in the orchestrator as a dict keyed by
\texttt{user\_id}, so two concurrent users cannot see each other's histories.

\subsection{The routing log}
\label{sec:log}

Every call to \texttt{Orchestrator.process} persists a row into
\texttt{routing\_log.db} with fields \texttt{(ts, user\_id, query, ml\_decision,
ml\_confidence, llm\_decision, llm\_confidence, final\_routing, model\_used)}
(\texttt{router/classification\_logger.py}). The logger exposes an
\texttt{export\_labeled\_csv(path, agreement\_only=True)} helper: by default
it exports only the rows where the ML and LLM agreed (a proxy for a correct
label), yielding a self-supervised training set for the next router
version. This closes the loop: more traffic $\rightarrow$ more labels
$\rightarrow$ better ML router $\rightarrow$ fewer LLM classifier calls.

\section{Experiments}
\label{sec:experiments}

We evaluate four properties of the system: (1) the offline router's accuracy
on a held-out split of the seed, (2) its qualitative behaviour on obvious
``small'' and ``big'' queries, (3) the correctness of the cascade under
mocked LLMs, and (4) the correctness of the conversation buffer under
arbitrary turn sequences.

\paragraph{Router accuracy.} We train on 108 hand-labelled queries with a
stratified 80/20 split and report test accuracy on the held-out 22 examples.
Seed labels were assigned by hand following the guideline: route ``small'' if
a competent 7B assistant should nail it, else ``big.''

\paragraph{Cascade integration tests.} We mock the small and large OpenAI
clients with a \texttt{FakeLLM} that records every call and can be scripted
per-prompt predicate. We then run the orchestrator through five end-to-end
scenarios: (a) high-confidence ``small'' skips the LLM classifier, (b)
high-confidence ``big'' routes to the large model and assembles intro/core/closing,
(c) low-confidence ML output cascades to the LLM classifier and its label
wins, (d) the pre-flight length check routes to the large model without
ever calling either router, (e) per-user buffer isolation.

\paragraph{Buffer correctness.} We exercise the ConversationBuffer with and
without a summarizer: empty-state invariants, single-turn verbatim
preservation, window-size eviction, degraded no-LLM mode, incremental
summary update across multiple evictions, assistant truncation, and a
schema check on \texttt{snapshot()}.

\paragraph{Setup.} All experiments run on a single MacBook Pro (Apple M-series
CPU, 16\,GB RAM) in a Python 3.14 virtualenv with scikit-learn 1.7, joblib
1.4, and pytest 8. Router training and inference are fully offline; the
integration tests never contact the OpenAI API (they use
\texttt{FakeLLM}). Total wall-clock time for \texttt{pytest -q} is under 5
seconds; router training is under 1 second on this seed set.

\section{Results}
\label{sec:results}

\begin{table}[h]
  \caption{Offline ML-router accuracy on the 22-example held-out split.
  Numbers averaged over a single stratified split (seed=42); see
  \S\ref{sec:limits} for why we do not report CIs.}
  \label{tab:router}
  \centering
  \begin{tabular}{lc}
    \toprule
    Metric & Value \\
    \midrule
    Test accuracy & $0.955$ \\
    Classes & $\{\mathrm{small}, \mathrm{big}\}$ \\
    Training examples & $86$ \\
    Test examples & $22$ \\
    Training wall-clock & $<1$ s \\
    \bottomrule
  \end{tabular}
\end{table}

\begin{table}[h]
  \caption{Pytest outcomes on the 38-test suite. All tests pass on a clean
  checkout after \texttt{python -m router.train}.}
  \label{tab:tests}
  \centering
  \begin{tabular}{llc}
    \toprule
    Module & Tests & Pass \\
    \midrule
    \texttt{test\_features.py}               & 10 & 10/10 \\
    \texttt{test\_conversation\_buffer.py}   & 9  & 9/9 \\
    \texttt{test\_ml\_router.py}             & 7  & 7/7 \\
    \texttt{test\_classification\_logger.py} & 3  & 3/3 \\
    \texttt{test\_orchestrator\_integration.py} & 9 & 9/9 \\
    \midrule
    Total                                    & 38 & 38/38 \\
    \bottomrule
  \end{tabular}
\end{table}

\paragraph{Qualitative routing behaviour.} On the four canonical ``big''
probes used in \texttt{test\_routes\_obvious\_big\_queries}
(\emph{Prove that there are infinitely many primes}, \emph{Design a
distributed rate limiter}, \emph{Compare Rust, Go, and Zig}, \emph{Derive
closed-form linear regression}), the router routes 4/4 to ``big''. On the
canonical ``small'' probes (\emph{Hi!}, \emph{What is the capital of
Japan?}, \emph{Thanks!}, \emph{What is 2 + 2?}), it routes $\geq 3/4$ to
``small.'' These are identical to the acceptance criteria encoded in the
test suite.

\paragraph{Cascade behaviour.} With a high-confidence small prediction the
orchestrator calls the small LLM exactly \emph{once} (the answer), skipping
the LLM classifier entirely. With a synthetic low-confidence ML prediction
(forced by monkeypatching \texttt{ml\_router.predict} to return
$c = 0.40$), the cascade calls the LLM classifier, which returns
\texttt{routing\_decision = big}; the final \texttt{routed\_by} field in the
response is \texttt{"llm"} and \texttt{ml\_prediction.confidence = 0.4},
exactly as specified. The pre-flight length bypass routes a 30-token input
to the large LLM with \texttt{routed\_by = "pre-flight"} and never touches
either router.

\paragraph{Buffer behaviour.} With \texttt{window\_size=1}, adding two
turn-pairs produces exactly one summarizer call; adding a third produces
two; the second summarizer call receives the previous summary (\texttt{SUM1})
in its user message, confirming that the summary is truly incremental and
not a re-summarization of the full trail. When the no-LLM fallback is used,
evicted turns still appear in \texttt{rolling\_summary} as tagged plain text,
which is enough for \texttt{build\_history()} to emit a valid history list.

\paragraph{Latency decomposition.} Because the small/big LLM calls dominate
wall-clock time, and because the ML router adds $\sim$1--2\,ms on our laptop,
the end-to-end effect of the cascade on the small route is
\emph{strictly one fewer round-trip}: a small-route turn that would
previously have required \texttt{classifier $\rightarrow$ answerer} now
requires just \texttt{answerer}, with the ML router running locally. On
the big route, the cascade saves the classifier round-trip only when the ML
router is confident; otherwise it adds the local ML-router cost on top of
the original LLM classifier call. Because the ML router is
sub-millisecond-cheap compared to a network round-trip, this is always a
net win in expectation.

\section{Discussion and Limitations}
\label{sec:limits}

\paragraph{Seed-set scale.} Our held-out test set is only 22 examples; a
95.5\% number on 22 examples has a Wilson 95\% lower bound near 78\%. We
therefore read the result as \emph{qualitatively strong on the axes we
targeted}, not as a tight generalization claim. Closing this gap is the
purpose of the routing log and its agreement-only CSV export: as real
traffic accumulates, the next router version can be retrained on thousands
rather than hundreds of labels.

\paragraph{Two-class labels.} Our router is deliberately binary. A
production system likely wants at least a third class for
``retrieval/tools'' and a fourth for ``refuse/deflect,'' which would require
a larger label set than our seed affords.

\paragraph{Summarization fidelity.} The conversation buffer uses the small
LLM to compress evicted turns. The small LLM can silently lose entities if a
rare proper noun is introduced and then not referenced again within the
window. Mitigation is straightforward---bump the window size, or preserve
named-entity tokens verbatim in the summary---but we did not evaluate it
formally here.

\paragraph{Profile leakage.} We deliberately keep the user profile out of the
large-LLM prompt, but the small LLM still sees it. If the small LLM runs on
the same API/vendor as the large LLM, the operational privacy benefit is
smaller than the architectural diagram suggests.

\paragraph{Threshold selection.} The cascade threshold $\tau = 0.65$ was
chosen by eye from the confidence histograms on the seed, not calibrated
against a held-out cost/quality trade-off. A proper treatment would jointly
optimize $\tau$ and the classifier under a cost function that prices small-
and big-route errors differently.

\section{Broader Impacts}
\label{sec:impacts}

\paragraph{Positive.} A fraction of every production chat workload is cheap
to serve. Systems that recognize this waste less energy per answered query,
lower the entry cost for student- and research-scale deployments, and
reduce latency for the simple queries that users make most frequently.

\paragraph{Negative.} A routing system that silently downgrades to a smaller
model can surprise users if quality varies visibly. Additionally, the
routing log is a per-query transcript and is therefore PII-sensitive;
deployments that enable the continuous-learning CSV export should treat the
log with the same retention, access, and consent controls as the underlying
conversation store. Our released code uses a local SQLite file with no
remote shipping and emits a warning on missing PII controls, but downstream
users must make deployment-specific decisions.

\section{Conclusion}
\label{sec:conclusion}

We showed that a $\sim$200-line offline classifier, trained on 108
hand-labelled examples and wired as the first stage of a cascade, reaches
95.5\% test accuracy on our seed and is strong enough in practice to skip
the LLM classifier on confident queries, cutting one API round-trip from
the common path. Paired with a bounded conversation buffer that
incrementally summarizes evicted turns, the resulting dual-LLM system
routes a long-running multi-user chat workload to the right model with
bounded per-turn cost. The full system (router, buffer, orchestrator, log,
and demo UI) is open-source and validated by a 38-test pytest suite.

\begin{ack}
We thank the instructors and teaching staff of our course at Northeastern
University for framing this as a final project with a clear
presentation/report split, which forced the system to be \emph{both}
demonstrable and documentable.

\textbf{Funding and competing interests.} This work was performed as
coursework; the authors received no external funding and declare no
competing interests.
\end{ack}

\section*{References}
\small

\begin{thebibliography}{9}

\bibitem{chen_frugalgpt_2023}
Chen, L., Zaharia, M., \& Zou, J. (2023).
FrugalGPT: How to use large language models while reducing cost and
improving performance. \emph{arXiv preprint arXiv:2305.05176}.

\bibitem{ong_routellm_2024}
Ong, I., Almahairi, A., Wu, V., Chiang, W.-L., Wu, T., Gonzalez, J.~E.,
Kadous, M.~W., \& Stoica, I. (2024).
RouteLLM: Learning to route LLMs with preference data.
\emph{arXiv preprint arXiv:2406.18665}.

\bibitem{databricks_routing_2024}
Databricks. (2024).
LLM Inference Performance Engineering: Model routing and cascades.
Technical blog post.

\bibitem{vaswani_attention_2017}
Vaswani, A., Shazeer, N., Parmar, N., Uszkoreit, J., Jones, L., Gomez, A.~N.,
Kaiser, {\L}., \& Polosukhin, I. (2017).
Attention is all you need. In \emph{Advances in Neural Information Processing
Systems 30}.

\bibitem{langchain_memory_2023}
LangChain. (2023).
ConversationSummaryBufferMemory --- LangChain documentation.
\url{https://python.langchain.com/docs/modules/memory/types/summary_buffer}.

\bibitem{schlag_rmt_2023}
Schlag, I., Irie, K., \& Schmidhuber, J. (2023).
Linear transformers are secretly fast weight programmers.
In \emph{International Conference on Machine Learning}.

\bibitem{zhang_personachat_2018}
Zhang, S., Dinan, E., Urbanek, J., Szlam, A., Kiela, D., \& Weston, J. (2018).
Personalizing dialogue agents: I have a dog, do you have pets too?
In \emph{ACL}.

\bibitem{salemi_lamp_2024}
Salemi, A., Mysore, S., Bendersky, M., \& Zamani, H. (2024).
LaMP: When large language models meet personalization.
In \emph{ACL}.

\bibitem{bommasani_foundation_2021}
Bommasani, R. et al. (2021).
On the opportunities and risks of foundation models.
\emph{arXiv preprint arXiv:2108.07258}.

\bibitem{openai_pricing_2024}
OpenAI. (2024).
Model pricing --- GPT-4o and GPT-4o-mini.
\url{https://openai.com/api/pricing/}.

\bibitem{our_repo}
Qi, Z., Sun, Y., \& Vijayarao, S. (2026).
dual-llm-system: source code.
\url{https://github.com/sandeepvijayarao09/dual-llm-system}.

\end{thebibliography}

\appendix

\section{Reproduction instructions}
\label{app:repro}

\begin{lstlisting}[language=bash]
git clone https://github.com/sandeepvijayarao09/dual-llm-system
cd dual-llm-system
python3 -m venv .venv && source .venv/bin/activate
pip install -r requirements.txt

# Train the v0 router (reproduces Table 1).
python -m router.train

# Run the full test suite (reproduces Table 2).
pytest -q

# Launch the demo UI (requires OPENAI_API_KEY in .env).
streamlit run app.py
\end{lstlisting}

\section{Key configuration defaults}

\begin{center}
\begin{tabular}{ll}
  \toprule
  Constant & Value \\
  \midrule
  Small LLM model                    & \texttt{gpt-4o-mini} \\
  Large LLM model                    & \texttt{gpt-4o} \\
  \texttt{ML\_ROUTER\_CONFIDENCE\_THRESHOLD} $\tau$ & $0.65$ \\
  \texttt{CONFIDENCE\_THRESHOLD}     & $0.70$ (LLM-classifier upgrade) \\
  \texttt{LARGE\_INPUT\_THRESHOLD}   & $200$ tokens \\
  \texttt{CONVERSATION\_WINDOW\_SIZE} $N$ & $3$ turn-pairs \\
  \texttt{max\_assistant\_chars} $M$ & $400$ chars \\
  TF-IDF \texttt{ngram\_range}       & $(1,2)$ \\
  LogisticRegression \texttt{C}      & $1.0$ \\
  Seed size                          & $108$ queries \\
  Train/test split                   & $80/20$ stratified \\
  \bottomrule
\end{tabular}
\end{center}

\section{Threats to validity}

Our seed set was authored by the same team that wrote the router, so label
bias correlates with implementation bias. The 22-example held-out split is
small and the split seed is fixed. The FakeLLM-based integration tests
validate control flow but not real model quality. All three caveats point at
the same next step: an online A/B evaluation against real traffic, with the
agreement-only log export feeding router v1.

\newpage
\section*{NeurIPS Paper Checklist}

\begin{enumerate}

\item {\bf Claims}
    \item[] Question: Do the main claims made in the abstract and introduction accurately reflect the paper's contributions and scope?
    \item[] Answer: \answerYes{}
    \item[] Justification: The five contributions bulleted in \S\ref{sec:intro} (cascade router, personalization-by-compression, bounded conversation buffer, routing log with CSV export, open-source implementation with a passing test suite) each have a corresponding section in the body (\S\ref{sec:router}, \S\ref{sec:compress}, \S\ref{sec:buffer}, \S\ref{sec:log}, \S\ref{sec:experiments}--\S\ref{sec:results}). We do not claim performance beyond the 22-example held-out split (Table~\ref{tab:router}).

\item {\bf Limitations}
    \item[] Question: Does the paper discuss the limitations of the work performed by the authors?
    \item[] Answer: \answerYes{}
    \item[] Justification: \S\ref{sec:limits} discusses seed-set scale and the weak Wilson lower bound at 22 examples, the binary-only label space, summarization fidelity risks, residual profile leakage to the small LLM, and the fact that $\tau$ was chosen by eye rather than calibrated.

\item {\bf Theory assumptions and proofs}
    \item[] Question: For each theoretical result, does the paper provide the full set of assumptions and a complete (and correct) proof?
    \item[] Answer: \answerNA{}
    \item[] Justification: The paper presents a system design and an empirical evaluation; it contains no theorems requiring proof. Equations~\ref{eq:cascade} and \ref{eq:summary} are definitions of the dispatch rule and incremental update, not claims.

\item {\bf Experimental result reproducibility}
    \item[] Question: Does the paper fully disclose all the information needed to reproduce the main experimental results of the paper to the extent that it affects the main claims and/or conclusions of the paper (regardless of whether the code and data are provided or not)?
    \item[] Answer: \answerYes{}
    \item[] Justification: Appendix~\ref{app:repro} gives the exact commands (\texttt{python -m router.train}, \texttt{pytest -q}); the seed data lives in \texttt{router/seed\_data.py} inside the released repo, and all hyperparameters (TF-IDF \texttt{ngram\_range}, logistic regression \texttt{C}, split seed, thresholds) are listed in Appendix~B.

\item {\bf Open access to data and code}
    \item[] Question: Does the paper provide open access to the data and code, with sufficient instructions to faithfully reproduce the main experimental results, as described in supplemental material?
    \item[] Answer: \answerYes{}
    \item[] Justification: Code, seed data, trained v0 model weights, and tests are publicly available at \url{https://github.com/sandeepvijayarao09/dual-llm-system} under an open-source license; Appendix~\ref{app:repro} gives a three-command reproduction path.

\item {\bf Experimental setting/details}
    \item[] Question: Does the paper specify all the training and test details (e.g., data splits, hyperparameters, how they were chosen, type of optimizer) necessary to understand the results?
    \item[] Answer: \answerYes{}
    \item[] Justification: \S\ref{sec:router} specifies the feature set and the full sklearn pipeline; \S\ref{sec:experiments} specifies the 80/20 stratified split; Appendix~B lists every threshold and hyperparameter used.

\item {\bf Experiment statistical significance}
    \item[] Question: Does the paper report error bars suitably and correctly defined or other appropriate information about the statistical significance of the experiments?
    \item[] Answer: \answerNo{}
    \item[] Justification: We report a single stratified split result rather than cross-validated error bars because the held-out set has only 22 examples; in \S\ref{sec:limits} we explicitly note the Wilson 95\% lower bound and flag this as the motivation for the routing-log-driven retraining loop.

\item {\bf Experiments compute resources}
    \item[] Question: For each experiment, does the paper provide sufficient information on the computer resources (type of compute workers, memory, time of execution) needed to reproduce the experiments?
    \item[] Answer: \answerYes{}
    \item[] Justification: \S\ref{sec:experiments} states that all experiments run on a single MacBook Pro (Apple M-series CPU, 16\,GB RAM), that \texttt{pytest -q} completes in under 5 seconds, and that router training completes in under 1 second.

\item {\bf Code of ethics}
    \item[] Question: Does the research conducted in the paper conform, in every respect, with the NeurIPS Code of Ethics \url{https://neurips.cc/public/EthicsGuidelines}?
    \item[] Answer: \answerYes{}
    \item[] Justification: The project uses only synthetic/hand-authored queries, no human-subject data, and public commercial APIs under their stated terms of service; the authors reviewed the NeurIPS Code of Ethics and found no conflicts.

\item {\bf Broader impacts}
    \item[] Question: Does the paper discuss both potential positive societal impacts and negative societal impacts of the work performed?
    \item[] Answer: \answerYes{}
    \item[] Justification: \S\ref{sec:impacts} discusses energy/cost reductions on the positive side and silent quality downgrades plus routing-log PII risk on the negative side.

\item {\bf Safeguards}
    \item[] Question: Does the paper describe safeguards that have been put in place for responsible release of data or models that have a high risk for misuse (e.g., pre-trained language models, image generators, or scraped datasets)?
    \item[] Answer: \answerNA{}
    \item[] Justification: The released router is a TF-IDF + logistic regression classifier trained on 108 benign hand-authored queries; it cannot generate text and poses no dual-use risk. The underlying LLMs (GPT-4o, GPT-4o-mini) are accessed via OpenAI's API under OpenAI's existing safeguards and are not released by us.

\item {\bf Licenses for existing assets}
    \item[] Question: Are the creators or original owners of assets (e.g., code, data, models), used in the paper, properly credited and are the license and terms of use explicitly mentioned and properly respected?
    \item[] Answer: \answerYes{}
    \item[] Justification: scikit-learn (BSD-3), joblib (BSD-3), pytest (MIT), Streamlit (Apache 2.0), python-dotenv (BSD-3), and the OpenAI Python SDK (Apache 2.0) are all used under their original licenses and cited in the references or the project \texttt{requirements.txt}. GPT-4o and GPT-4o-mini are used via the OpenAI API under OpenAI's API Terms of Use \citep{openai_pricing_2024}.

\item {\bf New assets}
    \item[] Question: Are new assets introduced in the paper well documented and is the documentation provided alongside the assets?
    \item[] Answer: \answerYes{}
    \item[] Justification: The released repository includes a \texttt{README.md}, per-module docstrings, typed dataclasses for the router's and orchestrator's return values, a trained \texttt{router\_v0.joblib} weights file, and a 38-test pytest suite that doubles as executable documentation.

\item {\bf Crowdsourcing and research with human subjects}
    \item[] Question: For crowdsourcing experiments and research with human subjects, does the paper include the full text of instructions given to participants and screenshots, if applicable, as well as details about compensation (if any)?
    \item[] Answer: \answerNA{}
    \item[] Justification: The paper involves no crowdsourcing and no human-subjects research; the 108-example seed was authored by the three authors themselves.

\item {\bf Institutional review board (IRB) approvals or equivalent for research with human subjects}
    \item[] Question: Does the paper describe potential risks incurred by study participants, whether such risks were disclosed to the subjects, and whether Institutional Review Board (IRB) approvals (or an equivalent approval/review based on the requirements of your country or institution) were obtained?
    \item[] Answer: \answerNA{}
    \item[] Justification: No human subjects, crowdworkers, or external data collection were involved, so IRB review does not apply.

\item {\bf Declaration of LLM usage}
    \item[] Question: Does the paper describe the usage of LLMs if it is an important, original, or non-standard component of the core methods in this research? Note that if the LLM is used only for writing, editing, or formatting purposes and does \emph{not} impact the core methodology, scientific rigor, or originality of the research, declaration is not required.
    \item[] Answer: \answerYes{}
    \item[] Justification: LLMs are a core component of the system under study: GPT-4o-mini is used as the small-route answerer, LLM classifier, summarizer, personalizer, and conversation-buffer compressor; GPT-4o is used as the big-route reasoner. Their roles are described in detail in \S\ref{sec:system}--\S\ref{sec:method} and summarized in Table~\ref{tab:components}. The authors additionally used an LLM-based coding assistant for code drafting and manuscript copyediting; it was not used for experimental design, analysis, or result generation.

\end{enumerate}


\end{document}
